\scanpage{39}
\mspage{38}

Ça montre, de plus, que si on part de $(M,{\leq},{\ll})$, il y a un très large choix pour les $\parallel$ possibles sur $\unsure{L|\mathcal L}$, i.e. les parties \unsure{envisagées} de $\unsure{\mathrm P|\mathcal P}_2(L)$. Si on désigne par
\[
 \unsure{\text{``}|\mathinner{..}}\quad
 \Gamma_0=\left\{\{x,y\}\in\mathrm P_2(L)\ \middle|\
 \begin{aligned}
 &\exists\,X,Y\in M,\ Y\leq X\\
 &y\deleted{\unsure{\leq}}\ll\deleted{\unsure{X}}\unsure{Y|X},\
 x\deleted{\unsure{\leq}}\ll X,\ x\not\ll Y
 \end{aligned}\right\}
\]
on peut prendre comme structure de graphe $\Gamma$ \deleted{\unsure{de}} (pour une relation $\parallel$) n'importe quelle partie $\Gamma\supset\Gamma_0$.
\[
 \Gamma_0\subset\Gamma\subset\mathrm P_2(L)\,.
\]

Ça fait de la marge ! Pour la suite, il faudra voir comment les autres axiomes affectent le choix de $\parallel$ (et du reste). Pour le moment, je songe surtout à
\[
 \begin{aligned}
 &\mathrm{At\ supp} &&\text{\unsure{(cf|cf} page 18)}\\
 &\mathrm{At\ \unsure{ens|ends}}&&\text{et ses variantes}\quad
 \text{(cf GF$\overline{\underline{\unsure{\mathrm{VI}|\mathrm{VII}}}}$ p 52 \unsure{....|...})}
 \end{aligned}
\]

\noindent\begin{tabular}{@{}||p{\dimexpr\linewidth-2\tabcolsep-3pt\relax}@{}}
\deleted{\unsure{Dans|Pour}} Pour qu'on ait un atlas ensembliste \inserted{spacieux}, \inserted{(resp. \unsure{polyédrisé|polyédralisé|polyédral})} suffit-il que $x\parallel y\Longleftrightarrow x\ne y$ ?? Sans
\end{tabular}
doute pas, car cela n'implique aucune restriction sur les autres données $\leq$, $\ll$, et les conditions $\mathrm{ML\,0}$ à $\mathrm{ML\,3}$ n'impliquent \uline{pas} que l'application \unsure{enlocus|unilocus|endroits|endroit|ombres} : $M\longrightarrow\mathrm{Fig}(L)$ soit injective, ni que $\ll$ soit induite par la relation $\ll$ entre fig. ensemblistes \inserted{dans $L$} \unsure{$\to$}
